\chapter{Attention, from one query to many heads}\label{ch:attention}

\begin{onething}
Attention is a similarity-weighted average: every token scores every other token with a dot
product, turns the scores into proportions that add up to one with a softmax, and replaces
itself with the matching blend of the other tokens' values.
\end{onething}

This chapter builds ordinary self-attention from nothing. \Cref{ch:esa} then shows the one
change SegFormer makes to it, and \cref{ch:block} wraps it into the repeating unit of the
encoder. Every formula here is followed by a calculation small enough to do on paper; please do
them on paper, with a pencil, because the arithmetic is the understanding. All the numbers in
this chapter were also checked by a small script, \file{research/textbook/tools/compute_13-15.py},
whose output is saved next to it.

\section{The problem: a token needs its surroundings}\label{sec:attention:problem}

After the patch embedding of \cref{ch:patchembed}, a $128\times128$ patch has become a grid of
$32\times32 = 1024$ \emph{tokens}. A token\index{token} is a short vector of numbers (32 of them at
stage~1) that summarises one $4\times4$ block of pixels. Recall from \cref{ch:shapes} that the
encoder holds all the tokens of a batch in a tensor of shape \shape{(B, N, C)}: $B$ patches in the
batch, $N$ tokens per patch, $C$ numbers per token.

So far each token knows only its own block. But the meaning of a block depends on what surrounds
it. A dark block in the middle of a lake is ordinary water. The identical dark block in the middle
of a bright desert is exactly the kind of thing an anomaly detector exists to notice. If the model
is to reconstruct a block from its context (the whole idea of \cref{ch:recon}), each token must be
able to look at other tokens and revise itself.

\paragraph{The obvious fixes, and why they fail.}
\begin{enumerate}
  \item \emph{Average the neighbours with fixed weights.} That is a convolution
  (\cref{ch:shapes}). Its weights are learned once and then applied identically everywhere,
  whatever the content. A convolution cannot decide that, for \emph{this} dark block, the
  relevant neighbours are the other dark blocks two rows away rather than the bright ones next
  door.
  \item \emph{Learn one weight for every pair of tokens.} With $N = 1024$ tokens that is
  $1024 \times 1024 = 1{,}048{,}576$ weights, and they would still be fixed: pair (3, 517) would
  get the same weight in a lake patch and in a desert patch.
\end{enumerate}
What we want is weights that are \emph{computed from the content}, freshly for every patch. That
is what attention provides.

\begin{analogy}
You walk into a library with a question. You do not read every book with equal care, and you do
not consult a fixed list of books regardless of the question. You read the titles on the spines
(keys), compare each with your question (query), and spend your time on the books whose titles
match best, taking from each a share of its contents (values) in proportion to how well it
matched. Change the question and you read a different blend of books. Every token in the encoder
is, at the same time, a reader with a question and a book on the shelf.
\end{analogy}

The mechanism has three ingredients: a way to measure similarity (the dot product,
\cref{sec:attention:dot}), a way to turn similarities into proportions (the softmax,
\cref{sec:attention:softmax}), and a way to blend (the weighted average,
\cref{sec:attention:onequery}). Then we assemble them in matrix form, explain the scaling by
$\sqrt{d}$, and split the work across several heads.

\section{Similarity: the dot product}\label{sec:attention:dot}

\begin{definition}[Dot product]
For two vectors $\vect{a} = (a_1, \dots, a_d)$ and $\vect{b} = (b_1, \dots, b_d)$ of the same
length $d$, the dot product\index{dot product} is
\begin{equation}\label{eq:attention:dot}
  \vect{a}\cdot\vect{b} = \vect{a}\T\vect{b} = \sum_{i=1}^{d} a_i b_i .
\end{equation}
\end{definition}

\Cref{ch:linalg} met it already; here is the property attention relies on. Writing
$\norm{\vect{a}} = \sqrt{\vect{a}\cdot\vect{a}}$ for the length of a vector and $\theta$ for the
angle between the two vectors,
\begin{equation}\label{eq:attention:cos}
  \vect{a}\cdot\vect{b} = \norm{\vect{a}}\,\norm{\vect{b}}\cos\theta .
\end{equation}
So the dot product is large and positive when the vectors point the same way ($\cos\theta$ near
1), near zero when they are at right angles, and negative when they point in opposite directions.
It is a similarity score, and it grows with the lengths of the vectors too, which will matter in
\cref{sec:attention:scale}.

\begin{worked}[label={we:attention:dots}]{Three dot products}
Take the query $\vect{q} = (1, 0)$ and three keys $\vect{k}_1 = (1, 0)$, $\vect{k}_2 = (0, 1)$,
$\vect{k}_3 = (1, 1)$. Multiply component by component and add.
\begin{align*}
  \vect{q}\cdot\vect{k}_1 &= 1\cdot 1 + 0\cdot 0 = 1 + 0 = 1,\\
  \vect{q}\cdot\vect{k}_2 &= 1\cdot 0 + 0\cdot 1 = 0 + 0 = 0,\\
  \vect{q}\cdot\vect{k}_3 &= 1\cdot 1 + 0\cdot 1 = 1 + 0 = 1.
\end{align*}
The query ``asks about'' the first feature. Keys 1 and 3 both have a first feature of 1, so
they score 1; key 2 has none, so it scores 0. Note that $\vect{k}_3$ is longer
($\norm{\vect{k}_3} = \sqrt{2}$) but the score depends only on the part of $\vect{k}_3$ that lies
along $\vect{q}$.
\end{worked}

\section{From scores to proportions: the softmax}\label{sec:attention:softmax}

\subsection{What we need, and the formula that delivers it}

Scores can be any real numbers, negative included. To take a weighted average we need weights
that are
\begin{enumerate}
  \item \emph{positive}, so that no token is subtracted;
  \item \emph{summing to one}, so that the result is an average and not a sum that grows with
  the number of tokens;
  \item \emph{order-preserving}, so that a higher score always gets a higher weight;
  \item \emph{smooth}, so that a small change in a score makes a small change in the weights
  (training needs derivatives, \cref{ch:trainloop}).
\end{enumerate}
Step 1, make every score positive while keeping the order: apply the exponential function
$e^{s}$, which is positive for every real $s$ and strictly increasing. Step 2, make them sum to
one: divide each by their total. The result is the softmax.

\begin{definition}[Softmax]
For scores $\vect{s} = (s_1, \dots, s_n)$,
\begin{equation}\label{eq:attention:softmax}
  \softmax(\vect{s})_i = \frac{e^{s_i}}{\sum_{j=1}^{n} e^{s_j}}, \qquad i = 1,\dots,n.
\end{equation}
\end{definition}\index{softmax}

Check the four requirements. Each numerator is positive and so is the denominator, so every
weight is positive. Adding all $n$ weights gives $\sum_i e^{s_i} / \sum_j e^{s_j} = 1$. Because
$e^{s}$ is increasing, $s_a > s_b$ gives $e^{s_a} > e^{s_b}$ and so a larger weight. And $e^s$ is
smooth.

The name says what it does: it is a ``soft'' version of picking the maximum. The largest score
gets the largest share, but every score gets something.

\begin{worked}[label={we:attention:softmax}]{Softmax of three scores}
Take $\vect{s} = (2, 1, 0)$. Do this on paper with a calculator for $e^x$.

\emph{Exponentiate each score.}
\begin{align*}
  e^{2} &= 7.3891, & e^{1} &= 2.7183, & e^{0} &= 1.0000 .
\end{align*}
\emph{Add them.}
\begin{align*}
  7.3891 + 2.7183 + 1.0000 &= 11.1074 .
\end{align*}
(To four decimals the exact total is 11.1073; the last digit differs because each term was
rounded first.)

\emph{Divide each by the total.}
\begin{align*}
  7.3891 / 11.1073 &= 0.6652,\\
  2.7183 / 11.1073 &= 0.2447,\\
  1.0000 / 11.1073 &= 0.0900.
\end{align*}
\emph{Check the sum.} $0.6652 + 0.2447 + 0.0900 = 0.9999$, which is 1 up to rounding.

The scores were evenly spaced (2, 1, 0), but the weights are not: each step down in score
divides the weight by $e \approx 2.718$. Softmax turns \emph{differences} of scores into
\emph{ratios} of weights.
\end{worked}

That last remark is worth making exact. For any two positions $a$ and $b$,
\begin{align}
  \frac{\softmax(\vect{s})_a}{\softmax(\vect{s})_b}
  &= \frac{e^{s_a} / \sum_j e^{s_j}}{e^{s_b} / \sum_j e^{s_j}}
  && \text{write out both weights} \notag\\
  &= \frac{e^{s_a}}{e^{s_b}} && \text{the common denominator cancels} \notag\\
  &= e^{s_a - s_b} && \text{rule } e^{x}/e^{y} = e^{x-y}. \label{eq:attention:ratio}
\end{align}
A score lead of 1 means a weight ratio of $e = 2.718$; a lead of 2 means $e^2 = 7.389$; a lead of
11 means about 60{,}000.

\subsection{Shifting all scores changes nothing}\label{sec:attention:shift}

\begin{keypoint}
Adding the same constant to every score leaves the softmax unchanged; only the differences
between scores matter.
\end{keypoint}

\emph{Proof.} Let $c$ be any number and add it to every score.
\begin{align*}
  \softmax(\vect{s} + c)_i
  &= \frac{e^{s_i + c}}{\sum_j e^{s_j + c}} && \text{definition, with } s_j + c \text{ in place of } s_j\\
  &= \frac{e^{s_i} e^{c}}{\sum_j e^{s_j} e^{c}} && \text{rule } e^{x+y} = e^{x} e^{y}\\
  &= \frac{e^{c}\, e^{s_i}}{e^{c} \sum_j e^{s_j}} && e^{c} \text{ does not depend on } j \text{, so it comes out of the sum}\\
  &= \frac{e^{s_i}}{\sum_j e^{s_j}} = \softmax(\vect{s})_i && \text{cancel } e^{c} \neq 0. \qquad\blacksquare
\end{align*}
This agrees with \cref{eq:attention:ratio}: the ratios depend only on differences.

\subsection{Numerical stability: subtract the maximum}\label{sec:attention:stable}

Computers store these numbers as 32-bit floating-point values (\emph{float32}). The largest
float32 is about $3.4028\times10^{38}$, and $\ln(3.4028\times10^{38}) = 88.72$. So $e^{s}$ for
any score above 88.72 cannot be stored; it becomes ``infinity''. In PyTorch,
\code{torch.exp(torch.tensor(89.0))} returns \code{inf}. Two infinities divided give
``not a number'', and one such value poisons a whole training run.

The cure follows from \cref{sec:attention:shift}. Choose $c = -\max_j s_j$: subtract the largest
score from all of them. The softmax is unchanged, but now the largest shifted score is exactly 0,
so the largest exponential is $e^0 = 1$ and nothing can overflow. Every other exponential lies
between 0 and 1. \index{softmax!numerical stability}

\begin{worked}[label={we:attention:stable}]{The stable softmax}
First redo \cref{we:attention:softmax} the stable way. The maximum of $(2, 1, 0)$ is 2.

\emph{Subtract the maximum.} $(2-2,\ 1-2,\ 0-2) = (0, -1, -2)$.

\emph{Exponentiate.} $e^{0} = 1.0000$, $e^{-1} = 0.3679$, $e^{-2} = 0.1353$.

\emph{Add.} $1.0000 + 0.3679 + 0.1353 = 1.5032$.

\emph{Divide.} $1.0000/1.5032 = 0.6652$, $0.3679/1.5032 = 0.2447$, $0.1353/1.5032 = 0.0900$.

The same weights as before, as the proof promised. Now a case where it matters: scores
$(1000, 999)$. Directly, $e^{1000}$ overflows. Stably:
\begin{align*}
  (1000 - 1000,\ 999 - 1000) &= (0, -1),\\
  e^{0} = 1.0000,\quad e^{-1} &= 0.3679,\\
  1.0000 + 0.3679 &= 1.3679,\\
  1.0000/1.3679 = 0.7311,\quad 0.3679/1.3679 &= 0.2689.
\end{align*}
PyTorch's \code{F.softmax}, which the model calls, does this subtraction internally.
\end{worked}

\subsection{Saturation: when one score takes everything}\label{sec:attention:saturation}

Stability stops the computer from failing. It does not stop softmax from becoming useless when
the scores are spread far apart. Consider scores $(20, 0, -15)$:
\begin{align*}
  e^{20} &= 485{,}165{,}195.4, & e^{0} &= 1, & e^{-15} &= 0.000\,000\,306 .
\end{align*}
The total is $485{,}165{,}196.4$ (the third term does not even reach the first decimal), so the
weights are
\begin{align*}
  485{,}165{,}195.4 / 485{,}165{,}196.4 &= 0.999\,999\,998,\\
  1 / 485{,}165{,}196.4 &= 0.000\,000\,002\,1,\\
  0.000\,000\,306 / 485{,}165{,}196.4 &= 0.000\,000\,000\,000\,000\,63 .
\end{align*}
The first token takes the whole average. We say the softmax has
\emph{saturated}\index{softmax!saturation}. Two things go wrong.
\begin{enumerate}
  \item The output is a copy of one value; no blending happens.
  \item Training cannot fix it, because, as we now show, the softmax's derivatives are
  essentially zero when it is saturated.
\end{enumerate}

\paragraph{The derivative of the softmax.} Training adjusts every number in the model in the
direction that lowers the loss (\cref{ch:recon} introduced the loss; \cref{ch:trainloop} does the
adjusting). To know that direction it needs derivatives. Write $p_i = \softmax(\vect{s})_i$ and
$Z = \sum_j e^{s_j}$, so $p_i = e^{s_i}/Z$. We want $\partial p_i / \partial s_k$, how much weight
$i$ changes when score $k$ moves.

First, $\partial Z / \partial s_k = e^{s_k}$, because only the $k$-th term of the sum contains
$s_k$, and the derivative of $e^{s}$ is $e^{s}$.

\emph{Case $k = i$.} Use the quotient rule $(u/v)' = (u'v - uv')/v^2$ with $u = e^{s_i}$,
$v = Z$:
\begin{align*}
  \frac{\partial p_i}{\partial s_i}
  &= \frac{e^{s_i} Z - e^{s_i} e^{s_i}}{Z^2} && \text{quotient rule}\\
  &= \frac{e^{s_i}}{Z}\cdot\frac{Z - e^{s_i}}{Z} && \text{take out the factor } e^{s_i}/Z\\
  &= p_i (1 - p_i) && \text{since } e^{s_i}/Z = p_i .
\end{align*}
\emph{Case $k \neq i$.} Now $u = e^{s_i}$ does not depend on $s_k$, so $u' = 0$:
\begin{align*}
  \frac{\partial p_i}{\partial s_k}
  &= \frac{0\cdot Z - e^{s_i} e^{s_k}}{Z^2} = -\frac{e^{s_i}}{Z}\cdot\frac{e^{s_k}}{Z} = -p_i p_k .
\end{align*}
Both cases in one line, with $\delta_{ik}$ equal to 1 when $i = k$ and 0 otherwise:
\begin{equation}\label{eq:attention:softmaxgrad}
  \frac{\partial p_i}{\partial s_k} = p_i(\delta_{ik} - p_k).
\end{equation}

Every derivative contains a product of two weights. If one weight is nearly 1 and the rest
nearly 0, every such product is nearly 0.

\begin{worked}[label={we:attention:satgrad}]{Gradients of a healthy and a saturated softmax}
Use $\partial p_i/\partial s_i = p_i(1-p_i)$.

\emph{Healthy scores $(1, 0, -1)$.} By the stable recipe the weights are $(0.6652, 0.2447,
0.0900)$ (the same as for $(2,1,0)$, by \cref{sec:attention:shift}). Then
\begin{align*}
  0.6652\times(1 - 0.6652) &= 0.6652\times0.3348 = 0.2227,\\
  0.2447\times(1 - 0.2447) &= 0.2447\times0.7553 = 0.1848,\\
  0.0900\times(1 - 0.0900) &= 0.0900\times0.9100 = 0.0819.
\end{align*}
\emph{Saturated scores $(20, 0, -15)$.} The weights are $(0.999\,999\,998,\ 2.1\times10^{-9},\
6.3\times10^{-16})$. Then
\begin{align*}
  0.999\,999\,998\times 0.000\,000\,002 &\approx 2.1\times10^{-9},\\
  2.1\times10^{-9}\times(1 - 2.1\times10^{-9}) &\approx 2.1\times10^{-9},\\
  6.3\times10^{-16}\times(1 - 6.3\times10^{-16}) &\approx 6.3\times10^{-16}.
\end{align*}
The healthy softmax passes about a fifth of any error signal back to its scores; the saturated
one passes about two parts in a billion. The losing tokens never get the signal that would let
them improve, and learning stalls.
\end{worked}

So we want scores whose spread is of order 1, not of order 20. \Cref{sec:attention:scale} shows
where large spreads come from and how one division prevents them.

\section{One query against many keys}\label{sec:attention:onequery}

We now have all three ingredients. Give one token a \emph{query} vector $\vect{q}$ and give every
token $j = 1,\dots,n$ a \emph{key} $\vect{k}_j$ and a \emph{value} $\vect{v}_j$, all of length
$d$. Then
\begin{align}
  s_j &= \vect{q}\cdot\vect{k}_j && \text{score: how relevant is token } j\text{?} \label{eq:attention:s}\\
  p_j &= \softmax(s_1,\dots,s_n)_j && \text{weight: its share of attention} \label{eq:attention:p}\\
  \vect{o} &= \sum_{j=1}^{n} p_j \vect{v}_j && \text{output: the weighted average of values.} \label{eq:attention:o}
\end{align}
(The full model also divides each score by $\sqrt{d}$; we add that in
\cref{sec:attention:scale}, and the worked example below already includes it so that its numbers
match the course.)

\begin{worked}[label={we:attention:onequery}]{One query, three tokens of dimension 2}
Three tokens, each with two features:
\[
  \vect{x}_1 = (1, 0),\qquad \vect{x}_2 = (0, 1),\qquad \vect{x}_3 = (1, 1).
\]
For now let the query, key and value of every token be the token itself (in the next section
they become learned transformations of it). Token 1 asks: $\vect{q} = \vect{x}_1 = (1, 0)$.

\emph{Step 1, scores.} From \cref{we:attention:dots}: $(1, 0, 1)$.

\emph{Step 1b, scale by $1/\sqrt{d}$ with $d = 2$.} $1/\sqrt{2} = 0.7071$, so the scaled scores
are $(0.7071,\ 0,\ 0.7071)$.

\emph{Step 2, softmax.} Exponentiate:
\begin{align*}
  e^{0.7071} &= 2.0281, & e^{0} &= 1.0000, & e^{0.7071} &= 2.0281.
\end{align*}
Add: $2.0281 + 1.0000 + 2.0281 = 5.0562$. Divide:
\begin{align*}
  2.0281/5.0562 &= 0.4011, & 1.0000/5.0562 &= 0.1978, & 2.0281/5.0562 &= 0.4011.
\end{align*}
Check: $0.4011 + 0.1978 + 0.4011 = 1.0000$.

\emph{Step 3, weighted average of the values,} one feature at a time.
First feature (the first number of each value):
\begin{align*}
  0.4011\times 1 &= 0.4011\\
  0.1978\times 0 &= 0.0000\\
  0.4011\times 1 &= 0.4011\\
  \text{sum} &= 0.8022 .
\end{align*}
Second feature:
\begin{align*}
  0.4011\times 0 &= 0.0000\\
  0.1978\times 1 &= 0.1978\\
  0.4011\times 1 &= 0.4011\\
  \text{sum} &= 0.5989 .
\end{align*}
So $\vect{o}_1 = (0.8022,\ 0.5989)$.

Token 1 went in as $(1, 0)$ and came out as $(0.80, 0.60)$. Its second feature was zero and is
now 0.60, because it judged tokens 2 and 3, which have a second feature, worth listening to.
It has absorbed context.
\end{worked}

\section{Queries, keys and values are learned projections}\label{sec:attention:qkv}

In \cref{we:attention:onequery} the query, key and value were all the token itself. That forces
a token to search for exactly what it advertises, and to hand over exactly what it is. Real
attention separates the three roles with three learned matrices.

\begin{definition}[Query, key, value]
Let $\vect{x}\in\R^{C}$ be a token. With three learned matrices
$\mat{W}_Q, \mat{W}_K, \mat{W}_V \in \R^{C\times C}$ and bias vectors
$\vect{b}_Q, \vect{b}_K, \vect{b}_V\in\R^{C}$,
\begin{align}
  \vect{q} = \vect{x}\T\mat{W}_Q + \vect{b}_Q\T,\qquad
  \vect{k} = \vect{x}\T\mat{W}_K + \vect{b}_K\T,\qquad
  \vect{v} = \vect{x}\T\mat{W}_V + \vect{b}_V\T .
\end{align}
The query\index{query} is what the token is looking for, the key\index{key} is what it
advertises about itself, and the value\index{value} is what it hands over when consulted.
\end{definition}

Each is a \emph{linear layer}\index{linear layer} (\cref{ch:shapes}): a matrix product plus a
bias. The learned numbers in $\mat{W}$ and $\vect{b}$ are \emph{parameters}: values the training
loop adjusts. In code (\file{app/foundation_models/components/efficient_self_attention.py},
lines 87--94):
\begin{lstlisting}
self.q = nn.Linear(embed_dim, embed_dim)     # (*, C) -> (*, C)
self.k = nn.Linear(embed_dim, embed_dim)
self.v = nn.Linear(embed_dim, embed_dim)
self.proj = nn.Linear(embed_dim, embed_dim)  # output projection, section 7
\end{lstlisting}
\begin{sourcenote}[: which way round PyTorch stores the matrix]
PyTorch's \code{nn.Linear(in, out)} stores its weight with shape \shape{(out, in)} and computes
$\vect{y} = \vect{x}\mat{W}\T + \vect{b}$, as the source comment at
\file{efficient_self_attention.py} line 89 says. This book writes $\vect{x}\T\mat{W}$ with
$\mat{W}$ of shape $(\text{in}, \text{out})$; the two are the same map with $\mat{W}$ stored
transposed. The parameter count, $\text{in}\times\text{out} + \text{out}$, is the same either way.
\end{sourcenote}

Why three separate matrices and not one? Because ``what I am looking for'' and ``what I
advertise'' are usually different. A token over a road edge might want to find other road
pixels (its query), while advertising itself as ``bright, linear'' (its key) and contributing
its exact spectrum-derived features (its value). Separate projections let a token search for one
thing while offering another.

\section{Attention in matrix form}\label{sec:attention:matrix}

Doing \cref{we:attention:onequery} once per token is slow to write and slow to compute. Matrices
do all tokens at once.

\subsection{Stacking tokens}

Stack the $N$ tokens as the rows of a matrix $\mat{X}$ of shape $(N, C)$; row $i$ is token $i$.
Then all queries at once are
\begin{equation}
  \mat{Q} = \mat{X}\mat{W}_Q,\qquad (N, C)\,(C, C) \to (N, C),
\end{equation}
because row $i$ of a product $\mat{X}\mat{W}$ is (row $i$ of $\mat{X}$) times $\mat{W}$. (We drop
the biases in this section to keep the algebra short; they add the same vector to every row.)
Likewise $\mat{K} = \mat{X}\mat{W}_K$ and $\mat{V} = \mat{X}\mat{W}_V$, each $(N, C)$.

\subsection{All scores at once: \texorpdfstring{$\mat{Q}\mat{K}\T$}{QK transpose}}

The transpose $\mat{K}\T$ swaps rows and columns, so it has shape $(C, N)$ and its column $j$ is
key $j$. Recall that entry $(i, j)$ of a product $\mat{A}\mat{B}$ is the dot product of row $i$ of
$\mat{A}$ with column $j$ of $\mat{B}$. Therefore
\begin{equation}\label{eq:attention:S}
  \mat{S} = \mat{Q}\mat{K}\T,\qquad (N, C)\,(C, N) \to (N, N),\qquad
  S_{ij} = \vect{q}_i\cdot\vect{k}_j .
\end{equation}
Entry $(i, j)$ is ``how relevant token $j$ is to token $i$''. Row $i$ holds all of token $i$'s
scores.

\subsection{The whole layer}

\begin{definition}[Scaled dot-product attention]
\begin{equation}\label{eq:attention:sdpa}
  \operatorname{Attention}(\mat{Q}, \mat{K}, \mat{V}) = \softmax_{\text{row}}\!\left(\frac{\mat{Q}\mat{K}\T}{\sqrt{d}}\right)\mat{V},
\end{equation}
where $\softmax_{\text{row}}$ applies \cref{eq:attention:softmax} to each row separately and $d$
is the length of each query and key.
\end{definition}\index{attention!scaled dot-product}
Shapes, step by step:
\begin{align*}
  \mat{Q}\mat{K}\T/\sqrt{d} &: (N, d)(d, N)\to(N, N) && \text{one score per (query, key) pair}\\
  \mat{A} = \softmax_{\text{row}}(\cdot) &: (N, N) && \text{each row sums to 1}\\
  \mat{A}\mat{V} &: (N, N)(N, d)\to(N, d) && \text{one blended value per token.}
\end{align*}
Row $i$ of $\mat{A}\mat{V}$ is $\sum_j A_{ij}\vect{v}_j$, exactly \cref{eq:attention:o} for query
$i$. The matrix $\mat{A}$ is the \emph{attention matrix}\index{attention matrix}.

This is the operation of \citet{vaswani2017attention}, and it is what the four lines of the
model's forward pass compute (\file{efficient_self_attention.py}, lines 182--195):
\begin{lstlisting}
scores  = Q_x @ K_x.transpose(-2, -1) * self.scale   # (B, h, N, N_red)
weights = F.softmax(scores, dim=-1)                   # rows sum to 1
out     = weights @ V_x                               # (B, h, N, d)
out     = self.proj(out_concatenated)                 # section 7
\end{lstlisting}
(\code{@} is matrix multiplication in Python; \code{dim=-1} means ``along the last axis'', that is,
along each row; \code{N_red} becomes $N$ until \cref{ch:esa}.)

\begin{worked}[label={we:attention:full}]{Full attention on three tokens, identity projections}
Use the tokens of \cref{we:attention:onequery} with $\mat{W}_Q=\mat{W}_K=\mat{W}_V=\mat{I}$, so
\[
  \mat{Q} = \mat{K} = \mat{V} = \mat{X} =
  \begin{pmatrix} 1 & 0\\ 0 & 1\\ 1 & 1 \end{pmatrix}.
\]
\emph{Transpose.} $\mat{K}\T = \begin{pmatrix} 1 & 0 & 1\\ 0 & 1 & 1\end{pmatrix}$.

\emph{Scores $\mat{Q}\mat{K}\T$,} one entry at a time (row of $\mat{Q}$ dotted with column of
$\mat{K}\T$):
\begin{align*}
  S_{11} &= 1\cdot1 + 0\cdot0 = 1, & S_{12} &= 1\cdot0 + 0\cdot1 = 0, & S_{13} &= 1\cdot1 + 0\cdot1 = 1,\\
  S_{21} &= 0\cdot1 + 1\cdot0 = 0, & S_{22} &= 0\cdot0 + 1\cdot1 = 1, & S_{23} &= 0\cdot1 + 1\cdot1 = 1,\\
  S_{31} &= 1\cdot1 + 1\cdot0 = 1, & S_{32} &= 1\cdot0 + 1\cdot1 = 1, & S_{33} &= 1\cdot1 + 1\cdot1 = 2.
\end{align*}
\emph{Scale} by $0.7071$:
\[
  \frac{\mat{S}}{\sqrt2} = \begin{pmatrix} 0.7071 & 0 & 0.7071\\ 0 & 0.7071 & 0.7071\\ 0.7071 & 0.7071 & 1.4142\end{pmatrix}.
\]
\emph{Softmax each row.} Row 1 was done in \cref{we:attention:onequery}: $(0.4011, 0.1978,
0.4011)$. Row 2 has the same numbers in a different order, so its weights are $(0.1978, 0.4011,
0.4011)$. Row 3 is new:
\begin{align*}
  e^{0.7071} = 2.0281,\quad e^{0.7071} &= 2.0281,\quad e^{1.4142} = 4.1133,\\
  2.0281 + 2.0281 + 4.1133 &= 8.1695,\\
  2.0281/8.1695 = 0.2483,\quad 2.0281/8.1695 &= 0.2483,\quad 4.1133/8.1695 = 0.5035.
\end{align*}
Check: $0.2483 + 0.2483 + 0.5035 = 1.0001$, one up to rounding. So
\[
  \mat{A} = \begin{pmatrix} 0.4011 & 0.1978 & 0.4011\\ 0.1978 & 0.4011 & 0.4011\\ 0.2483 & 0.2483 & 0.5035\end{pmatrix}.
\]
\emph{Multiply by $\mat{V}$.} Row 1 gave $(0.8022, 0.5989)$. Row 2:
\begin{align*}
  \text{feature 1: } 0.1978\cdot1 + 0.4011\cdot0 + 0.4011\cdot1 &= 0.1978 + 0 + 0.4011 = 0.5989,\\
  \text{feature 2: } 0.1978\cdot0 + 0.4011\cdot1 + 0.4011\cdot1 &= 0 + 0.4011 + 0.4011 = 0.8022.
\end{align*}
Row 3:
\begin{align*}
  \text{feature 1: } 0.2483\cdot1 + 0.2483\cdot0 + 0.5035\cdot1 &= 0.2483 + 0 + 0.5035 = 0.7518,\\
  \text{feature 2: } 0.2483\cdot0 + 0.2483\cdot1 + 0.5035\cdot1 &= 0 + 0.2483 + 0.5035 = 0.7518.
\end{align*}
(The script gives 0.7517 from unrounded weights.) The output is
\[
  \mat{A}\mat{V} = \begin{pmatrix} 0.8022 & 0.5989\\ 0.5989 & 0.8022\\ 0.7518 & 0.7518\end{pmatrix}.
\]
Tokens 1 and 2 are mirror images, as the symmetry of the input demands, and token 3, which
resembles both, has become an even blend.
\end{worked}

\begin{figure}[htbp]
\centering
\begin{tikzpicture}[x=1.45cm,y=1.0cm]
  % cells: row/col/value (rows top to bottom)
  \foreach \r/\c/\v in {1/1/0.4011,1/2/0.1978,1/3/0.4011,
                         2/1/0.1978,2/2/0.4011,2/3/0.4011,
                         3/1/0.2483,3/2/0.2483,3/3/0.5035}{
    \pgfmathtruncatemacro{\p}{\v*140}
    \fill[attn!\p] (\c-1, -\r+1) rectangle (\c, -\r);
    \draw[white, line width=1pt] (\c-1, -\r+1) rectangle (\c, -\r);
    \pgfmathsetmacro{\tc}{\v > 0.45 ? 1 : 0}
    \ifdim\tc pt>0.5pt
      \node[white, font=\small\bfseries] at (\c-0.5, -\r+0.5) {\v};
    \else
      \node[black, font=\small] at (\c-0.5, -\r+0.5) {\v};
    \fi
  }
  \foreach \c/\t in {1/{key 1},2/{key 2},3/{key 3}}
    \node[font=\small, spat] at (\c-0.5, 0.35) {\t};
  \foreach \r/\t in {1/{query 1},2/{query 2},3/{query 3}}
    \node[font=\small, spat, anchor=east] at (-0.1, -\r+0.5) {\t};
  \foreach \r/\s in {1/{0.4011 + 0.1978 + 0.4011 = 1},2/{0.1978 + 0.4011 + 0.4011 = 1},3/{0.2483 + 0.2483 + 0.5035 = 1}}
    \node[font=\footnotesize, anchor=west] at (3.15, -\r+0.5) {row sum: $\s$};
  \draw[decorate,decoration={brace,amplitude=4pt,mirror}] (0,-3.15) -- (3,-3.15)
    node[midway, below=5pt, font=\small] {softmax runs along each row};
\end{tikzpicture}
\caption{The attention matrix $\mat{A}$ of \cref{we:attention:full}, drawn as a heat map (darker
plum = larger weight). Row $i$ is how query $i$ shares its attention among the three keys; every
row sums to 1, columns need not. Token 3, which resembles both others, puts half its attention on
itself.}
\label{fig:attention:heatmap}
\end{figure}

\coursefig[0.95\linewidth]{iw-12-attention}{The same example as the course draws it: (a) three
2-number tokens with $\mat{Q}=\mat{K}=\mat{V}=\mat{X}$; (b) scaled scores $\vect{q}\cdot\vect{k}/\sqrt2$;
(c) softmax along each row; (d) each output is the weighted average of the values. Every number
matches \cref{we:attention:full}.}{fig:attention:course}

\begin{worked}[label={we:attention:proj}]{Full attention with real projections}
Now let the three roles differ. Keep $\mat{X}$ as before and take
\[
  \mat{W}_Q = \begin{pmatrix}1 & 0\\ 1 & 1\end{pmatrix},\qquad
  \mat{W}_K = \mat{I},\qquad
  \mat{W}_V = \begin{pmatrix}0 & 1\\ 1 & 0\end{pmatrix}.
\]
$\mat{W}_V$ swaps the two features: what a token hands over is its features in reverse order.

\emph{Queries $\mat{Q} = \mat{X}\mat{W}_Q$,} row by row (row of $\mat{X}$ times each column of
$\mat{W}_Q$):
\begin{align*}
  (1,0)\mat{W}_Q &= (1\cdot1 + 0\cdot1,\ 1\cdot0 + 0\cdot1) = (1, 0),\\
  (0,1)\mat{W}_Q &= (0\cdot1 + 1\cdot1,\ 0\cdot0 + 1\cdot1) = (1, 1),\\
  (1,1)\mat{W}_Q &= (1\cdot1 + 1\cdot1,\ 1\cdot0 + 1\cdot1) = (2, 1).
\end{align*}
\emph{Keys} $\mat{K} = \mat{X}$. \emph{Values $\mat{V} = \mat{X}\mat{W}_V$:} $(0,1)$, $(1,0)$,
$(1,1)$.

\emph{Scores $\mat{Q}\mat{K}\T$:}
\begin{align*}
  \text{row 1, } \vect{q}=(1,0)&: \quad 1,\ 0,\ 1\\
  \text{row 2, } \vect{q}=(1,1)&: \quad 1\cdot1+1\cdot0 = 1,\ \ 1\cdot0+1\cdot1 = 1,\ \ 1\cdot1+1\cdot1 = 2\\
  \text{row 3, } \vect{q}=(2,1)&: \quad 2\cdot1+1\cdot0 = 2,\ \ 2\cdot0+1\cdot1 = 1,\ \ 2\cdot1+1\cdot1 = 3 .
\end{align*}
\emph{Scale by 0.7071 and softmax.} Rows 1 and 2 are rows we have met: $(0.4011, 0.1978,
0.4011)$ and $(0.2483, 0.2483, 0.5035)$. Row 3, scaled, is $(1.4142, 0.7071, 2.1213)$:
\begin{align*}
  e^{1.4142} = 4.1133,\quad e^{0.7071} &= 2.0281,\quad e^{2.1213} = 8.3421,\\
  4.1133 + 2.0281 + 8.3421 &= 14.4835,\\
  4.1133/14.4835 = 0.2840,\quad 2.0281/14.4835 &= 0.1400,\quad 8.3421/14.4835 = 0.5760.
\end{align*}
\emph{Blend the values.} Row 1: $0.4011(0,1) + 0.1978(1,0) + 0.4011(1,1)$:
\begin{align*}
  \text{feature 1: } 0 + 0.1978 + 0.4011 &= 0.5989,\qquad
  \text{feature 2: } 0.4011 + 0 + 0.4011 = 0.8022 .
\end{align*}
Row 2: $0.2483(0,1) + 0.2483(1,0) + 0.5035(1,1)$ gives $(0.2483 + 0.5035,\ 0.2483 + 0.5035) =
(0.7518, 0.7518)$. Row 3: $0.2840(0,1) + 0.1400(1,0) + 0.5760(1,1)$:
\begin{align*}
  \text{feature 1: } 0 + 0.1400 + 0.5760 &= 0.7160,\qquad
  \text{feature 2: } 0.2840 + 0 + 0.5760 = 0.8600 .
\end{align*}
Compare with \cref{we:attention:full}. Token 2 used to attend mostly to itself and token 3;
with its new query $(1,1)$ it now prefers token 3, the one that has both features. The projections
changed \emph{who} is listened to (through $\mat{W}_Q$) and \emph{what} is heard (through
$\mat{W}_V$). Training learns these matrices.
\end{worked}

\section[Why divide by the square root of d]{Why divide by \texorpdfstring{$\sqrt{d}$}{the square root of d}}\label{sec:attention:scale}

\Cref{sec:attention:saturation} showed that softmax saturates when the scores are spread over
tens of units. Where would such spreads come from? From the length of the vectors.

\subsection{The variance of a dot product}

Recall from \cref{ch:linalg} that the \emph{variance}\index{variance} of a random number $X$ is
$\operatorname{Var}(X) = \E[X^2] - (\E[X])^2$, the average squared distance from its mean; its
square root is the \emph{standard deviation}, the typical size of a deviation.

\paragraph{Assumptions.} Model the components of a query and a key as random numbers:
$q_1,\dots,q_d$ and $k_1,\dots,k_d$, all independent, each with mean 0 and variance 1. (This is roughly
what LayerNorm, \cref{ch:block}, and standard weight initialisation produce at the start of
training.)

\emph{Step 1: one term $q_i k_i$ has mean 0.} For independent numbers the mean of a product is the
product of the means:
\[
  \E[q_i k_i] = \E[q_i]\,\E[k_i] = 0\cdot0 = 0 .
\]
\emph{Step 2: one term has variance 1.} Using the definition of variance and then independence
again (now for $q_i^2$ and $k_i^2$):
\begin{align*}
  \operatorname{Var}(q_i k_i) &= \E[(q_i k_i)^2] - (\E[q_i k_i])^2 && \text{definition}\\
  &= \E[q_i^2 k_i^2] - 0^2 && \text{step 1}\\
  &= \E[q_i^2]\,\E[k_i^2] && \text{independence}\\
  &= 1\cdot1 = 1 && \E[q_i^2] = \operatorname{Var}(q_i) + (\E[q_i])^2 = 1 + 0 .
\end{align*}
\emph{Step 3: add $d$ independent terms.} The variance of a sum of independent random numbers is
the sum of their variances. The terms $q_1k_1, \dots, q_dk_d$ are independent of each other, so
\begin{equation}\label{eq:attention:vard}
  \operatorname{Var}(\vect{q}\cdot\vect{k}) = \sum_{i=1}^{d}\operatorname{Var}(q_i k_i) = \underbrace{1 + 1 + \dots + 1}_{d\text{ terms}} = d .
\end{equation}
The standard deviation of a raw score is therefore $\sqrt{d}$: it grows with the head dimension.

\emph{Step 4: divide by $\sqrt d$.} For any constant $a$, $\operatorname{Var}(aX) =
a^2\operatorname{Var}(X)$ (every deviation is multiplied by $a$, every squared deviation by $a^2$).
With $a = 1/\sqrt d$:
\begin{equation}\label{eq:attention:var1}
  \operatorname{Var}\!\left(\frac{\vect{q}\cdot\vect{k}}{\sqrt d}\right) = \frac{1}{d}\cdot d = 1 .
\end{equation}
Whatever the head dimension, the scaled scores have a standard deviation of about 1, which keeps
the softmax in its gradual, trainable range.\index{attention!scaling by root d}

\begin{worked}[label={we:attention:enum}]{The variance of a 2-dimensional dot product, exactly}
Let each of $q_1, q_2, k_1, k_2$ be $+1$ or $-1$ with equal probability. Each has mean
$\tfrac12(1) + \tfrac12(-1) = 0$ and variance $\tfrac12(1)^2 + \tfrac12(-1)^2 = 1$, so the
assumptions hold. There are $2^4 = 16$ equally likely combinations.

Each product $q_ik_i$ is $+1$ (when the signs agree) or $-1$ (when they differ), each with
probability $\tfrac12$. The score $q_1k_1 + q_2k_2$ is therefore
\begin{align*}
  (+1) + (+1) &= +2 && \text{probability } \tfrac12\cdot\tfrac12 = \tfrac14 \quad(4\text{ of the 16})\\
  (+1) + (-1) \text{ or } (-1) + (+1) &= 0 && \text{probability } \tfrac14 + \tfrac14 = \tfrac12 \quad(8\text{ of the 16})\\
  (-1) + (-1) &= -2 && \text{probability } \tfrac14 \quad(4\text{ of the 16}).
\end{align*}
The mean is $\tfrac14(2) + \tfrac12(0) + \tfrac14(-2) = 0$. The variance is the mean of the
squares:
\begin{align*}
  \tfrac14(2)^2 + \tfrac12(0)^2 + \tfrac14(-2)^2 &= \tfrac14\cdot4 + 0 + \tfrac14\cdot4 = 1 + 0 + 1 = 2 = d .
\end{align*}
Dividing every score by $\sqrt2$ gives scores $\pm1.4142$ and 0, with variance $2/2 = 1$.
\end{worked}

The script repeats the experiment with 200{,}000 random Gaussian queries and keys for larger
$d$\tagC:

\begin{center}
\begin{tabular}{rrrrr}
\toprule
$d$ & sample variance of $\vect{q}\cdot\vect{k}$ & sample std & $\sqrt d$ & variance after $\div\sqrt d$\\
\midrule
2 & 2.002 & 1.415 & 1.414 & 1.001\\
16 & 15.985 & 3.998 & 4.000 & 0.999\\
32 & 31.746 & 5.634 & 5.657 & 0.992\\
64 & 64.029 & 8.002 & 8.000 & 1.000\\
\bottomrule
\end{tabular}
\end{center}

\paragraph{What this means in this model.} The head dimension at stages 2--4 is 32 (next
section). Unscaled, a typical score would be $\pm5.66$, and a score two standard deviations above
another ($2\times5.657 = 11.31$ higher) would get $e^{11.31} \approx 81{,}937$ times the weight, by
\cref{eq:attention:ratio}. Scaled, the same two-standard-deviation lead is 2, a weight ratio of
$e^2 = 7.39$. The first is a single-token copy; the second is a blend. The code
(\file{efficient_self_attention.py}, line 82) stores the factor once:
\begin{lstlisting}
self.scale = self.head_dim ** -0.5     # 1 / sqrt(d); d = 16 gives 0.25
\end{lstlisting}
(\code{** -0.5} is ``to the power minus one half'', which is one over the square root.)

\section{Attention ignores order: permutation equivariance}\label{sec:attention:perm}

Here is a property that surprises most readers. Shuffle the tokens before attention, and the
outputs come out shuffled in exactly the same way, with the same values. Attention has no idea
where any token sits.

\emph{Why.} Let $\mat{P}$ be a \emph{permutation matrix}: the identity matrix with its rows
reordered, so that $\mat{P}\mat{X}$ is $\mat{X}$ with its rows reordered. Feed $\mat{P}\mat{X}$
to attention.
\begin{align*}
  \mat{Q}' &= \mat{P}\mat{X}\mat{W}_Q = \mat{P}\mat{Q},\quad \mat{K}' = \mat{P}\mat{K},\quad \mat{V}' = \mat{P}\mat{V} && \text{rows reordered}\\
\intertext{The scores use the rule for the transpose of a product, $(\mat{A}\mat{B})\T = \mat{B}\T\mat{A}\T$ (for example, a $1\times2$ row times a $2\times1$ column is a number, and transposing swaps the order of the factors):}
  \mat{Q}'\mat{K}'^{\top} &= \mat{P}\mat{Q}(\mat{P}\mat{K})\T = \mat{P}\mat{Q}\mat{K}\T\mat{P}\T = \mat{P}\mat{S}\mat{P}\T\\
\intertext{The softmax works row by row, and reordering the rows of $\mat{S}$ and the entries within each row only reorders the softmax's output the same way:}
  \softmax_{\text{row}}(\mat{P}\mat{S}\mat{P}\T) &= \mat{P}\,\softmax_{\text{row}}(\mat{S})\,\mat{P}\T = \mat{P}\mat{A}\mat{P}\T\\
  \mat{A}'\mat{V}' &= \mat{P}\mat{A}\mat{P}\T\mat{P}\mat{V} = \mat{P}\mat{A}\mat{V} && \mat{P}\T\mat{P} = \mat{I}.
\end{align*}
The output is the old output with its rows reordered by the same $\mat{P}$. A function with this
property is called \emph{permutation equivariant}\index{permutation equivariance}.

\begin{worked}[label={we:attention:perm}]{Swapping two tokens}
Swap tokens 1 and 2 of \cref{we:attention:full}: the input becomes $(0,1), (1,0), (1,1)$. Redo
row 1: query $(0,1)$ against keys $(0,1), (1,0), (1,1)$ scores $1, 0, 1$, exactly the old row 1
of $\mat{S}$, so the weights are $(0.4011, 0.1978, 0.4011)$ and the output is
\begin{align*}
  0.4011(0,1) + 0.1978(1,0) + 0.4011(1,1) &= (0 + 0.1978 + 0.4011,\ 0.4011 + 0 + 0.4011)\\
  &= (0.5989,\ 0.8022),
\end{align*}
which is the old output of token 2. The script confirms the whole output is the old one with rows
1 and 2 swapped: $(0.5989, 0.8022), (0.8022, 0.5989), (0.7517, 0.7517)$\tagC.
\end{worked}

For a picture this is a problem: ``next to'' and ``far from'' are among the most useful facts
about a token, and attention cannot see them. Most transformers add a learned \emph{positional
encoding}\index{positional encoding} vector to every token to say where it is. This codebase has
none: a search of \file{app/} for \code{pos_embed} finds nothing. SegFormer gets its sense of
position from a $3\times3$ depthwise convolution inside the feed-forward network
\citep{xie2021segformer}, which \cref{ch:block} explains.

\section{Many heads}\label{sec:attention:heads}

\subsection{Split, attend, concatenate, project}

One attention gives one blend per token: a single compromise between everything the token might
want to know. Multi-head attention runs $h$ smaller attentions side by side, each on its own slice
of the features, so that different heads can learn different notions of relevance (one might
attend to near neighbours, another to anything of similar brightness anywhere in the patch).

\begin{definition}[Multi-head attention]
Let $C$ be divisible by $h$ and let $d = C/h$, the \emph{head dimension}\index{head dimension}.
Compute $\mat{Q}, \mat{K}, \mat{V}$ of shape $(N, C)$ as before and cut each into $h$ blocks of
$d$ columns: $\mat{Q} = [\mat{Q}^{(1)}\,|\,\cdots\,|\,\mat{Q}^{(h)}]$, and likewise for
$\mat{K}$ and $\mat{V}$. Then
\begin{align}
  \mat{O}^{(m)} &= \softmax_{\text{row}}\!\left(\frac{\mat{Q}^{(m)}\mat{K}^{(m)\top}}{\sqrt d}\right)\mat{V}^{(m)}, \qquad m = 1,\dots,h, \label{eq:attention:head}\\
  \operatorname{MHA}(\mat{X}) &= [\mat{O}^{(1)}\,|\,\cdots\,|\,\mat{O}^{(h)}]\,\mat{W}_O + \vect{b}_O\T . \label{eq:attention:mha}
\end{align}
Each head has its own $(N, N)$ attention matrix. The outputs, each $(N, d)$, are concatenated
back to $(N, C)$ and passed through an output projection\index{output projection}
$\mat{W}_O\in\R^{C\times C}$.
\end{definition}\index{attention!multi-head}

The output projection matters. Without it the $h$ heads would never interact; the model would be
$h$ separate narrow models stapled side by side. $\mat{W}_O$ lets every output feature draw on
every head.

\begin{analogy}
The library again. One reader forms one impression. Five specialists reading the same shelf (a
lawyer, an accountant, an engineer, a historian, an editor) each notice different things; the
output projection is the meeting where their notes are merged into one report.
\end{analogy}

\subsection{The reshaping, with shapes}

With a batch, $\mat{Q}$ is a tensor \shape{(B, N, C)}. The code splits the last axis into
(head, feature within head) and moves the head axis forward so each head's data is a contiguous
\shape{(N, d)} matrix (\file{efficient_self_attention.py}, lines 128--134 and 190--195):
\begin{lstlisting}
Q_x = self.q(x)                                        # (B, N, C)
Q_x = torch.reshape(Q_x, (B, N, self.num_heads, self.head_dim))  # (B, N, h, d)
Q_x = Q_x.transpose(1, 2)                              # (B, h, N, d)
# ... K_x, V_x likewise; attention runs on the last two axes for every (b, head)
out = out.transpose(1, 2)                              # (B, N, h, d)
out = torch.reshape(out, (B, N, self.head_dim * self.num_heads))  # (B, N, C)
out = self.proj(out)                                   # (B, N, C)
\end{lstlisting}
Matrix multiplication with \code{@} on 4-axis tensors treats the first two axes as a batch: for
every patch $b$ and head $m$ it multiplies the \shape{(N, d)} and \shape{(d, N)} matrices on the
last two axes. So \shape{(B, h, N, d) @ (B, h, d, N)} gives \shape{(B, h, N, N)}.

\begin{worked}[label={we:attention:split}]{Splitting a 4-number token into two heads}
Let $C = 4$, $h = 2$, so $d = 4/2 = 2$. A token's query $(0.5, -1, 2, 0)$ is cut into head 1's
query $(0.5, -1)$ and head 2's query $(2, 0)$. If a key is $(1, 1, 1, -3)$, head 1 scores
$0.5\cdot1 + (-1)\cdot1 = -0.5$ and head 2 scores $2\cdot1 + 0\cdot(-3) = 2$, each then divided by
$\sqrt2$. The full 4-number dot product would be $-0.5 + 2 = 1.5$: one head sees the key as
irrelevant, the other as highly relevant, and a single head would have averaged the two views
away.
\end{worked}

\subsection{The heads in this model}

\begin{center}
\begin{tabular}{crrrc}
\toprule
stage & $C$ (\code{embed_dim}) & $h$ (\code{num_heads}) & $d = C/h$ & whole number?\\
\midrule
1 & 32 & 2 & $32/2 = 16$ & yes\\
2 & 64 & 2 & $64/2 = 32$ & yes\\
3 & 160 & 5 & $160/5 = 32$ & yes\\
4 & 256 & 8 & $256/8 = 32$ & yes\\
\bottomrule
\end{tabular}
\end{center}
Widths and heads are those of the shipped fine-tune FT04: \code{embed_dims}
$= [32, 64, 160, 256]$ and \code{num_heads} $= [2, 2, 5, 8]$
\par\noindent\tagM{research/experiments/indradhanu-finetune/FT04-final-joint/config.json}.
\par\noindent Version 0.1.0 used $[1, 2, 5, 8]$\tagM{configs/hyperspectral_segformer_exp_1.json},
which gave stage 1 a head dimension of 32.

\begin{sourcenote}[: the code's default head count is not the trained one]
The constructors \file{app/foundation_models/components/hyperspectral_seg_former_mae.py}
(line 71) and \file{spectral_stem_seg_former_mae.py} (line 54) default to
\code{num_heads = [1, 2, 5, 8]}, and the encoder's docstring table
(\file{seg_former_encoder.py}, line 78) lists 1 head at stage 1. Every Indradhanu run since
v0.2.0 passes $[2, 2, 5, 8]$ explicitly (\file{configs/hyperspectral_segformer_exp_2.json}
line 33, the FT04 config above, and the backend's
\file{backend/allotrope/foundation_models/resolver.py} line 489). The course table follows the
configs and is right for the trained models; a model built with no arguments would have one head
at stage 1. As the next subsection shows, the parameter count is the same either way.
\end{sourcenote}

\subsection{Parameter count: \texorpdfstring{$4C^2 + 4C$}{4C squared plus 4C}}

The learned parts of plain multi-head attention are four linear layers, each $C\to C$: the matrices
$\mat{W}_Q$, $\mat{W}_K$, $\mat{W}_V$ and $\mat{W}_O$, and their biases. A linear layer from
$C_\text{in}$ to $C_\text{out}$ has $C_\text{in}C_\text{out}$ weights and $C_\text{out}$ biases
(\cref{ch:shapes}). So
\begin{align}
  \text{one projection} &= C\cdot C + C = C^2 + C, \notag\\
  \text{four projections} &= 4(C^2 + C) = 4C^2 + 4C . \label{eq:attention:params}
\end{align}
\emph{The number of heads does not appear.} Splitting into heads reshapes the outputs of the same
four matrices; it adds no weights. The script builds the stage-1 attention module with one head
and with two and gets 69{,}856 parameters both times\tagC\ (that total includes the
spatial-reduction convolution of \cref{ch:esa}).

\begin{worked}[label={we:attention:params}]{The projection parameters at every stage}
Apply \cref{eq:attention:params} stage by stage, squaring first, then multiplying, then adding.
\begin{align*}
  C = 32:&\quad 32^2 = 1{,}024;\ \ 4\times1{,}024 = 4{,}096;\ \ 4\times32 = 128;\ \ 4{,}096 + 128 = 4{,}224\\
  C = 64:&\quad 64^2 = 4{,}096;\ \ 4\times4{,}096 = 16{,}384;\ \ 4\times64 = 256;\ \ 16{,}384 + 256 = 16{,}640\\
  C = 160:&\quad 160^2 = 25{,}600;\ \ 4\times25{,}600 = 102{,}400;\ \ 4\times160 = 640;\\
  &\quad 102{,}400 + 640 = 103{,}040\\
  C = 256:&\quad 256^2 = 65{,}536;\ \ 4\times65{,}536 = 262{,}144;\ \ 4\times256 = 1{,}024;\\
  &\quad 262{,}144 + 1{,}024 = 263{,}168
\end{align*}
\emph{Check against the model.} Building the real \code{EfficientSelfAttention} modules, the
script reports \code{q}, \code{k}, \code{v}, \code{proj} at 1{,}056 each for $C=32$
($4\times1{,}056 = 4{,}224$), 4{,}160 each for $C=64$, 25{,}760 each for $C=160$, and 65{,}792 each
for $C=256$\tagC. At stage 4, where there is no reduction convolution, the whole attention module
is 263{,}168, exactly our number.
\end{worked}

\subsection{Cost}

Per head, $\mat{Q}\mat{K}\T$ takes $N\cdot N\cdot d$ multiply-adds (each of $N^2$ entries is a
$d$-term dot product) and $\mat{A}\mat{V}$ another $N\cdot N\cdot d$. Over $h$ heads that is
$2hN^2d = 2N^2C$: the work grows with the \emph{square} of the number of tokens.

\section{Self-attention and cross-attention}\label{sec:attention:cross}

Everything in this chapter is \emph{self}-attention\index{attention!self}: the queries, keys and
values all come from the same set of tokens. In \emph{cross}-attention\index{attention!cross} the
queries come from somewhere else (a small set of learned ``askers'') while the keys and values
still come from the tokens. The four steps are identical; only the source of $\mat{Q}$ changes,
and with it the shape of the score matrix, which becomes (number of askers)$\times N$. Part IV of
this book uses exactly that to compress a spectrum (\cref{ch:ceiling,ch:crossstem}).

\coursefig[0.9\linewidth]{iw-32-self-cross}{Where the questions come from. (a) In self-attention
the tokens produce queries, keys and values. (b) In cross-attention an outside set of queries
reads the tokens' keys and values. Right: the two score matrices for the spectral case of
\cref{ch:ceiling}, drawn to scale.}{fig:attention:selfcross}

\section{The cost problem}\label{sec:attention:cost}

The score matrix has one entry per (query, key) pair: $N\times N$ per head, per patch, per block.
For a $128\times128$ patch:
\begin{center}
\begin{tabular}{crrr}
\toprule
stage & grid & $N$ & $N\times N$\\
\midrule
1 & $32\times32$ & 1024 & $1024\times1024 = 1{,}048{,}576$\\
2 & $16\times16$ & 256 & $256\times256 = 65{,}536$\\
3 & $8\times8$ & 64 & $64\times64 = 4{,}096$\\
4 & $4\times4$ & 16 & $16\times16 = 256$\\
\bottomrule
\end{tabular}
\end{center}
Stage 1 has over a million scores for one head of one block of one patch: 4{,}096 times stage 4.
And it is the stage where every other tensor is largest. \Cref{ch:esa} is the fix.

\begin{confusions}
\confusion{Is attention the same as correlation?}{Related but not identical. Correlation is a
fixed statistic of two signals. Attention compares learned projections of the tokens, so training
decides what ``similar'' means for this task.}
\confusion{Do the rows or the columns of the attention matrix sum to one?}{The rows. Softmax runs
along the last axis (\texttt{dim=-1}), which is the key axis, so each query distributes exactly
one unit of attention. A column sum says how popular a key is and can be anything from 0 to $N$
(in \cref{fig:attention:heatmap} column 3 sums to 1.3057).}
\confusion{Why three separate matrices instead of comparing tokens directly?}{Because what a token
looks for, what it advertises, and what it hands over are different things. \Cref{we:attention:proj}
shows one projection changing who is listened to and another changing what is heard.}
\confusion{Do more heads mean more parameters?}{No. The head count only decides how the $C$
features are cut up; the four $C\times C$ matrices are the same size. It changes the head
dimension $d = C/h$, the scale $1/\sqrt d$, and how many separate attention patterns are formed.}
\confusion{Does the scale $1/\sqrt d$ change which token wins?}{No. Dividing every score by the
same positive number keeps their order. It changes how decisively the winner wins: the gaps
between scores shrink, so the weights are less extreme.}
\confusion{What does ``self'' mean in self-attention?}{That queries, keys and values all come
from the same token set. Cross-attention takes its queries from elsewhere (\cref{sec:attention:cross}).}
\end{confusions}

\begin{checkyourself}
\item\label{ex:attention:1} Compute $\softmax(3, 1, 0)$ by hand, first directly and then by the
stable recipe (subtract the maximum). Show every exponential, the sum and each division, and check
that the weights sum to 1.
\item\label{ex:attention:2} Redo \cref{we:attention:onequery} for token 2, $\vect{q} = (0, 1)$,
with identity projections. Give the scores, scaled scores, exponentials, weights and output.
\item\label{ex:attention:3} Show that for two scores the softmax weight of the first is
$1/(1 + e^{-(s_1 - s_2)})$, the logistic sigmoid of the score difference. What weight does a lead
of $s_1 - s_2 = 2$ give?
\item\label{ex:attention:4} Suppose the query and key components are independent with mean 0 but
variance $\sigma^2 = 4$ instead of 1. Rederive $\operatorname{Var}(\vect{q}\cdot\vect{k})$ step by
step for head dimension $d = 16$. What divisor would bring the variance back to 1?
\item\label{ex:attention:5} Stage 3 has $C = 160$ and $h = 5$. (a) What is the head dimension and
the scale factor? (b) Count the parameters of its four projections by hand. (c) The whole stage-3
attention module has 205{,}920 parameters; how many belong to something other than the four
projections?
\item\label{ex:attention:6} Stage 4 has $C = 256$, $h = 8$ and $N = 16$. Give the shapes of
$\mat{Q}$ after the head split, of the score tensor for a batch of $B$ patches, and of the output
after concatenation. How many score entries are there per patch, summed over all heads?
\end{checkyourself}

\nextchapter{ch:esa}
